%% ED26-P07-monotonicity-0001. Assigned source: Icard--Moss--Tune 2017.
\chapter{Monotonicity and Typed Inference}
\label{sec:p07-monotonicity-inference}

Replacing an expression by a more inclusive one can preserve an inference,
reverse its direction, or invalidate it. The answer depends on the position
occupied by the expression. This entry develops the marked typed calculus
of Icard, Moss and Tune \cite{IcardMossTune2017Monotonicity}, with the
interpretation of type coercions made explicit. Its conclusions are
inequalities in specified ordered domains. Natural-language uses additionally
require the lexical and compositional premises stated in the application.
\index{monotonicity!typed inference}\index{natural logic!polarity}
\index{antitone function}\index{polarity!occurrence}

\section{Types, markings and coercions}

\begin{hbdefinition}{def:p07-monotonicity-types}{Marked types and their order}
Fix a set \(B\) of base types. The markings are
\(\mathsf{Mar}=\{+,-,\cdot\}\), ordered by
\( +\sqsubseteq\cdot\), \( -\sqsubseteq\cdot\), and reflexivity.
Here \(+\) asserts monotonicity, \(-\) antitonicity, and \(\cdot\) makes
neither assertion. A function bearing the last marking can happen to be
monotone. The composition of markings is multiplication of signs on
\(\{+,-\}\), and is \(\cdot\) if either factor is \(\cdot\).
Thus \(+\) is its identity and the operation is associative.
Indeed, if one of three factors is \(\cdot\), both groupings are
\(\cdot\); otherwise this is multiplication of the numbers \(1,-1\).

The set \(T\) is the least set containing \(B\) and
\(\sigma\xrightarrow{m}\tau\) whenever \(\sigma,\tau\in T\) and
\(m\in\mathsf{Mar}\). Arrows associate to the right.
The height of a base type is \(0\); the height of an arrow is one plus
the maximum of the heights of its argument and result types.
The type preorder \(\preceq\) is the least preorder satisfying
\[
 \sigma'\preceq\sigma,\quad \tau\preceq\tau',\quad
 m\sqsubseteq m'
 \quad\Longrightarrow\quad
 (\sigma\xrightarrow m\tau)
 \preceq(\sigma'\xrightarrow{m'}\tau').
\]
This is an order of admissible coercions, contravariant in the argument type.
It is not a formation rule giving a term additional intrinsic types.

The source's auxiliary compatibility relation \(\uparrow\) and partial
operation \(\vee\) are specified by
\(\sigma\uparrow\sigma\) and \(\sigma\vee\sigma=\sigma\).
Whenever \(\tau_1\uparrow\tau_2\), it also declares
\((\sigma\xrightarrow{m_1}\tau_1)\uparrow
  (\sigma\xrightarrow{m_2}\tau_2)\)
and defines their join to be
\(\sigma\xrightarrow{m_1\vee m_2}(\tau_1\vee\tau_2)\).
For distinct markings \(m_1,m_2\), \(m_1\vee m_2=\cdot\).
No value is assigned when these clauses do not define one.
Finally,
\(\widehat b=b\) and
\(\widehat{\sigma\xrightarrow m\tau}=
  \sigma\xrightarrow{\cdot}\widehat\tau\).
This last operation changes markings along the result spine; it does not
erase all markings in argument types.
\end{hbdefinition}

\begin{hblemma}{thm:p07-monotonicity-type-order}
 {proved=annex:p7.monotonicity.type-order}{Structural characterization of the type order}
Distinct base types are incomparable. A base type and an arrow type are
incomparable. For arrow types,
\[
 \sigma\xrightarrow m\tau\preceq\sigma'\xrightarrow{m'}\tau'
 \quad\Longleftrightarrow\quad
 \sigma'\preceq\sigma,\quad \tau\preceq\tau',\quad m\sqsubseteq m'.
\]
In particular, the relation is antisymmetric. If the auxiliary
\(\sigma\vee\tau\) is defined, both inputs are below it.
Not every pair of types has such an auxiliary join.
\end{hblemma}

\begin{hbdefinition}{def:p07-monotonicity-prestructure}{Coherent prestructures}
For every \(\tau\in T\), let \(D_\tau\) be a set with a reflexive,
transitive relation \(\leq_\tau\). For every \(\sigma\preceq\tau\), supply
an order-preserving map \(\pi_{\sigma,\tau}:D_\sigma\to D_\tau\).
Require identity and composition:
\[
 \pi_{\sigma,\sigma}=\operatorname{id},\qquad
 \pi_{\sigma,\upsilon}=\pi_{\tau,\upsilon}\pi_{\sigma,\tau}
 \quad(\sigma\preceq\tau\preceq\upsilon).
\]
At arrow types require actual functions, with
\[
 D_{\sigma\xrightarrow+\tau}\cup D_{\sigma\xrightarrow-\tau}
 \subseteq D_{\sigma\xrightarrow{\cdot}\tau}
 \subseteq \{f:D_\sigma\to D_\tau\}.
\]
The \(+\) functions preserve order; the \(-\) functions reverse order.
Every arrow domain carries pointwise order. Coercions at arrows act by
precomposition and postcomposition:
\[
 \pi_{\sigma\xrightarrow m\tau,\,
       \sigma'\xrightarrow{m'}\tau'}(f)(a)
 =
 \pi_{\tau,\tau'}\bigl(f(\pi_{\sigma',\sigma}(a))\bigr).
 \tag{C}
\]
Here the source type is below the target and \(a\in D_{\sigma'}\).
Forgetting only a marking is consequently the inclusion of the same
function. Equation (C) is part of the semantic specification throughout
this entry.
\end{hbdefinition}

\begin{hbexample}{def:p07-monotonicity-full-domains}{Full and restricted function domains}
Given base preorders, recursively take all functions at a \(\cdot\)-arrow,
all monotone functions at a \(+\)-arrow and all antitone functions at a
\(-\)-arrow. Use pointwise orders and (C). These are full prestructures.
The structural characterization above verifies the construction recursively:
composition with monotone coercions preserves a required \(+\) or \(-\)
marking, and pointwise order makes every coercion monotone.

Restricted function domains are also allowed, provided they satisfy the
same conditions and contain the denotations of all abstractions.
For \(B=\{r\}\) and \(D_r=\mathbb R\) with its usual order,
\(D_{r\xrightarrow-\!r}\) contains \(a\mapsto 1-a\), and
\(D_{r\xrightarrow+\!r}\) contains \(a\mapsto a-1\).
\end{hbexample}

\section{Terms and occurrence polarity}

\begin{hbdefinition}{def:p07-monotonicity-terms}{Intrinsic types and marked occurrences}
Choose constants, each with one assigned type, and infinitely many
variables of each type. A variable or constant is a term of its assigned
type. If \(M:\sigma\xrightarrow m\tau\) and \(N:\rho\preceq\sigma\),
then \(M(N):\tau\). The intrinsic type is denoted \(\operatorname{tp}(M)\).

Each free occurrence of a variable in itself has polarity \(+\).
An occurrence in the function part \(M\) of \(M(N)\) retains its polarity.
An occurrence of polarity \(n\) in the argument \(N\) acquires polarity
\(m\circ n\). Constants have no free occurrences.

For \(x:\sigma\) and \(M:\tau\), the term \(\lambda x.M\) has type
\(\sigma\xrightarrow m\tau\), where \(m=+\) if there is at least one free
\(x\) in \(M\) and all such occurrences are positive; \(m=-\) if there is
at least one and all are negative; otherwise \(m=\cdot\).
In particular, vacuous abstraction has intrinsic marking \(\cdot\),
although its denotation is both monotone and antitone.
Occurrences of other variables retain their polarities under abstraction.
Terms are identified up to renaming bound variables by fresh variables of
the same type. Every term has exactly one intrinsic type.
\end{hbdefinition}

\begin{hbdefinition}{def:p07-monotonicity-structures}{Structures and interpretation}
An assignment \(\varphi\) gives \(x:\sigma\) a value in \(D_\sigma\).
Interpret each constant by a fixed element of its assigned domain.
An interpretation on a coherent prestructure satisfies
\[
 \sem{x}_\varphi=\varphi(x),\qquad
 \sem{c}_\varphi=c^D,\qquad
 \sem{M(N)}_\varphi
 =\sem{M}_\varphi\bigl(\pi_{\rho,\sigma}(\sem N_\varphi)\bigr),
\]
where \(M:\sigma\xrightarrow m\tau\) and \(N:\rho\preceq\sigma\).
Writing \(\varphi_x^a\) for the assignment changed only at \(x\), require
\[
 \sem{\lambda x.M}_\varphi(a)=\sem M_{\varphi_x^a}.
\]
Every value belongs to the term's intrinsic domain. A structure is a
prestructure equipped with such an interpretation; closure under all
abstractions is therefore required. Whenever a term \(M:\rho\) is compared
at a supertype \(\tau\), its value there is
\(\pi_{\rho,\tau}(\sem M_\varphi)\).

In full prestructures these clauses define an interpretation recursively.
The following result verifies that the abstraction clauses have the required
markings. It also applies to restricted structures.
\end{hbdefinition}

\begin{hblemma}{thm:p07-monotonicity-polarity}
 {proved=annex:p7.monotonicity.polarity}{Occurrence polarity controls variation}
Fix a structure, an assignment, a term \(M:\tau\), and a variable
\(x:\sigma\). If every free occurrence of \(x\) in \(M\) is positive,
then \(a\mapsto\sem M_{\varphi_x^a}\) is monotone.
If every such occurrence is negative, this function is antitone.
When \(x\) is absent, both assertions apply and the function is constant.
\end{hblemma}

\begin{hbexample}{def:p07-monotonicity-real-terms}{Real-function examples}
Take \(0,1:r\), addition of type
\(r\xrightarrow+(r\xrightarrow+r)\), and subtraction of type
\(r\xrightarrow+(r\xrightarrow-r)\), with their usual meanings.
Then
\[
 \begin{array}{c|c|c}
 \text{term}&\text{intrinsic type}&\text{denotation}\\ \hline
 +(0)(1)&r&1\\
 +(1)(1)&r&2\\
 -(0)(1)&r&-1\\
 x&r&\varphi(x)\\
 -(x)&r\xrightarrow-r&b\mapsto\varphi(x)-b\\
 -(0)&r\xrightarrow-r&b\mapsto-b\\
 -(1)(x)&r&1-\varphi(x)\\
 -(x)(1)&r&\varphi(x)-1\\
 -(x)(y)&r&\varphi(x)-\varphi(y)\\
 \lambda x.-(1)(x)&r\xrightarrow-r&a\mapsto1-a\\
 \lambda x.-(x)(1)&r\xrightarrow+r&a\mapsto a-1
 \end{array}
\]
If \(f:r\xrightarrow{\cdot}r\), then
\(\lambda f.\lambda x.-(0)(f(x))\) has intrinsic type
\((r\xrightarrow{\cdot}r)\xrightarrow{-}
 (r\xrightarrow{\cdot}r)\).
Its denotation is pointwise negation of a function. The outer negative
marking follows because the free occurrence of \(f\) is in the negative
second argument of subtraction. The inner \(\cdot\) marking records that
the occurrence of \(x\) passes through the unspecified argument position
of \(f\). No monotonicity of an arbitrary input function is assumed.
\end{hbexample}

\section{Substitution, conversion and finite assertions}

\begin{hbdefinition}{def:p07-monotonicity-substitution}{Typed substitution and conversion}
A substitution \(s\) assigns each \(x:\sigma\) a term \(s(x)\) of type
\(\preceq\sigma\). Substitution is simultaneous and capture avoiding:
constants remain fixed, variables are replaced by their assigned terms,
and application is treated componentwise. Before descending beneath a
binder, rename it to a variable of the same type fresh for the finitely
many substitution images relevant to the body. Thus
\[
 (\lambda x.M)[s]=\lambda y.(M[s_x^y])
\]
with \(y\) fresh. Here \(s_x^y\) replaces the assignment to \(x\) by \(y\).
Write \(M[N/x]\) for a one-variable substitution.

The relation \(\approx\) is the congruence generated by alpha conversion,
capture-avoiding beta conversion
\((\lambda x.M)N\approx M[N/x]\), and eta conversion
\(\lambda x.M(x)\approx M\), with \(x\notin\FV(M)\).
Every displayed term and application must be well typed. Eta conversion
allows the binder type to be below the intrinsic argument type of \(M\).
It compares the function with its corresponding restriction; semantic
comparisons always name the type at which that restriction is used.
Congruence includes abstraction and application, with their actual
well-typedness conditions. A beta-eta normal term has neither kind of
redex. Mark-erasure gives an ordinary simply typed lambda term.
\end{hbdefinition}

\begin{hblemma}{thm:p07-monotonicity-substitution}
 {proved=annex:p7.monotonicity.substitution}{Substitution, type refinement and interpretation}
For a term \(M:\tau\), \(M[s]\) is a term of a type
\(\tau_s\preceq\tau\). If an occurrence of \(x\) has polarity \(m\) in
\(M\), and an occurrence of \(y\) has polarity \(n\) in \(s(x)\), every
corresponding surviving occurrence in \(M[s]\) has a polarity
\(\sqsubseteq m\circ n\).
In particular, a positive or negative polarity remains that same sign
whenever the product is a sign.

Put
\[
 \varphi_s(x)=
 \pi_{\operatorname{tp}(s(x)),\operatorname{tp}(x)}
       (\sem{s(x)}_\varphi).
\]
Then
\[
 \pi_{\tau_s,\tau}(\sem{M[s]}_\varphi)=\sem M_{\varphi_s}.
 \tag{S}
\]
Beta reduction refines types and preserves every tracked positive or
negative free-occurrence polarity. Both sign-refinement and the
possibility that substitution erases or duplicates occurrences are included.
\end{hblemma}

\begin{hblemma}{thm:p07-monotonicity-conversion}
 {proved=annex:p7.monotonicity.conversion}{Conversion is stable under substitution}
If \(M\approx N\), then \(M[s]\approx N[s]\) for every typed,
capture-avoiding substitution \(s\).
Beta reduction terminates by erasure into ordinary simply typed terms.
No uniqueness assertion about the remaining marked binder types is needed.
The ordinary normalization dependency is
Theorem~\ref{thm:p07-stlc-strong-normalization}, with inert symbols covered by
Corollary~\ref{thm:p07-stlc-inert-constants}.
\end{hblemma}

\begin{hbdefinition}{def:p07-monotonicity-conversion-coherence}
 {The conversion-coherence property}
A structure has conversion coherence when, whenever \(M\approx N\)
and both intrinsic types
are below \(\tau\), at every assignment
\[
 \pi_{\operatorname{tp}(M),\tau}(\sem M_\varphi)
 =
 \pi_{\operatorname{tp}(N),\tau}(\sem N_\varphi).
 \tag{E}
\]
Use an explicitly specified common upper type; an undefined auxiliary
join is never required. The property will follow from comparison-directed
elaboration, rather than being imposed as an additional model axiom.
\end{hbdefinition}

\begin{hblemma}{thm:p07-monotonicity-ordinary-long}
 {proved=annex:p7.monotonicity.ordinary-long}{The ordinary long-form input}
In the arrow-only ordinary simply typed lambda calculus, with a fixed
typed context and inert constants, every term has a finite beta-normal,
eta-long representative. Two terms of the same type are beta-eta
convertible exactly when these representatives are alpha-equivalent.
Here long means that every function-typed result is abstracted over its
arguments and every neutral head is fully applied to a base-typed result;
the arguments are recursively long.
This is the unmarked normalization input used below
\cite[Sections 4.2--4.4 and 8.7]{Selinger2013LambdaNotesV2}.
\end{hblemma}

\begin{hbdefinition}{def:p07-monotonicity-long-form}
 {Comparison-directed long forms}
For \(\operatorname{tp}(M)\preceq\tau\), let \(L_\tau(M)\) be the
beta-normal, eta-long elaboration at expected type \(\tau\).
Use fresh bound variables of the expected argument types.
At an expected arrow \(\tau=\sigma\xrightarrow m\rho\), form
\(\lambda u:\sigma.\,L_\rho(M(u))\), beta-normalizing the application
before descending. At an expected base type, a beta-normal term is a
variable or constant head with an application spine; elaborate each
argument at the corresponding argument type declared by that head.
Existing abstractions are processed by replacing their binders by the
fresh expected-type variables. Recompute intrinsic markings from the
resulting occurrences. The following lemma establishes that this
procedure is finite, well typed and independent of the choices of fresh
names up to typed alpha conversion.
\end{hbdefinition}

\begin{hblemma}{thm:p07-monotonicity-long-coherence}
 {proved=annex:p7.monotonicity.long-coherence}
 {Long elaboration and semantic coherence}
For every \(\operatorname{tp}(M)\preceq\tau\):
\begin{enumerate}
\item \(L_\tau(M)\) is a term of intrinsic type below \(\tau\), and
      \(L_\tau(M)\approx M\).
      Every tracked signed free occurrence survives with the same sign,
      possibly in several copies; an unspecified occurrence can refine.
\item In every coherent structure,
\[
 \pi_{\operatorname{tp}(L_\tau(M)),\tau}
        (\sem{L_\tau(M)}_\varphi)
 =
 \pi_{\operatorname{tp}(M),\tau}(\sem M_\varphi).
\]
\item If \(M\approx N\) and both are comparable at \(\tau\), their
      long elaborations at \(\tau\) are typed alpha-equivalent.
\end{enumerate}
Consequently every coherent structure has property (E).
Uniqueness here concerns the elaboration at a fixed comparison type,
not raw marked short normal forms.
\end{hblemma}

\begin{hbdefinition}{def:p07-monotonicity-consequence}{Assertions and universal premises}
An assertion is \(M\precsim_\tau N\), where
\(\operatorname{tp}(M),\operatorname{tp}(N)\preceq\tau\).
It is true at an assignment precisely when
\[
 \pi_{\operatorname{tp}(M),\tau}(\sem M_\varphi)
 \leq_\tau
 \pi_{\operatorname{tp}(N),\tau}(\sem N_\varphi).
\]
A structure satisfies an assertion when it is true at every assignment.
For a set \(\Gamma\) of assertions, write
\(\Gamma\models_{\mathrm{mon}}M\precsim_\tau N\) when every structure
satisfying every member of \(\Gamma\) satisfies the conclusion.
Thus free variables in premises are universally interpreted.
A derivation \(\Gamma\vdash_{\mathrm{mon}}M\precsim_\tau N\) is a finite
tree built from members of \(\Gamma\) and the rules on
Sheet~\ref{sheet:p07-monotonicity}.
\end{hbdefinition}

\showcalculus{p07-monotonicity}

\begin{hblemma}{thm:p07-monotonicity-instantiation}
 {proved=annex:p7.monotonicity.instantiation}{Instantiation and extensional comparison}
The following hold in the calculus.
\begin{enumerate}
 \item If \(M\approx M'\), \(N\approx N'\), and all four terms are
       comparable at \(\tau\), then
       \(\Gamma\vdash M\precsim_\tau N\) implies
       \(\Gamma\vdash M'\precsim_\tau N'\).
 \item For every typed substitution \(s\),
       \[
       M\precsim_\tau N\ \vdash_{\mathrm{mon}}\
       M[s]\precsim_\tau N[s].
       \]
       Consequently derivability from \(\Gamma\) is closed under
       substitution of its conclusion while retaining \(\Gamma\).
 \item Suppose \(M,N\) have intrinsic types below
       \(\sigma\xrightarrow m\tau\). If
       \(\Gamma\vdash M(O)\precsim_\tau N(O)\) for every term
       \(O:\sigma\), then
       \(\Gamma\vdash M\precsim_{\sigma\xrightarrow m\tau}N\).
\end{enumerate}
\end{hblemma}

\section{Representing syntax by ordered functions}

\begin{hbdefinition}{def:p07-monotonicity-term-structure}{Term structures}
Write \(\langle M\rangle\) for a conversion class and set
\[
 T_\tau=\{\langle M\rangle:
            \operatorname{tp}(M)\preceq\tau\}.
\]
This means that the class has such a representative. Application of
classes is computed with representatives that type-check; congruence
makes the result independent of that choice.

A term structure supplies each \(T_\tau\) with a preorder
\(\sqsubseteq_\tau\). A class in \(T_{\sigma\xrightarrow+\tau}\)
acts monotonically on \(T_\sigma\); a class in
\(T_{\sigma\xrightarrow-\tau}\) acts antitonically; and
\[
 \langle M\rangle\sqsubseteq_{\sigma\xrightarrow m\tau}\langle N\rangle
 \quad\Longleftrightarrow\quad
 \langle M(O)\rangle\sqsubseteq_\tau\langle N(O)\rangle
 \text{ for every }\langle O\rangle\in T_\sigma .
\]
The inclusions \(i_{\sigma,\tau}:T_\sigma\to T_\tau\) compose and have
identity components.
\end{hbdefinition}

\begin{hblemma}{thm:p07-monotonicity-term-inclusions}
 {proved=annex:p7.monotonicity.term-inclusions}{Order preservation between term domains}
The inclusion maps of a term structure preserve its preorders.
\end{hblemma}

\begin{hblemma}{thm:p07-monotonicity-term-model}
 {proved=annex:p7.monotonicity.term-model}
 {Comparison-indexed functional realization}
Every term structure has a realization by a coherent structure \(D\)
and order isomorphisms \(\iota_\tau:T_\tau\to D_\tau\).
The realization depends on its comparison type. For legal representatives,
\[
 \iota_{\sigma\xrightarrow m\tau}(\langle M\rangle)
       \bigl(\iota_\sigma(\langle O\rangle)\bigr)
   =\iota_\tau(\langle M(O)\rangle),
\qquad
 \pi_{\sigma,\tau}
   =\iota_\tau i_{\sigma,\tau}\iota_\sigma^{-1}.
\]
Every valuation \(\varphi\) is represented by a typed substitution \(s\),
and the interpretation is
\[
 \sem M_\varphi
    =\iota_{\operatorname{tp}(M)}(\langle M[s]\rangle).
\]
The displayed index is the original type of \(M\), even when substitution
refines its intrinsic type. The construction uses choice to select
representatives for a whole valuation.
\end{hblemma}

\begin{hbtheorem}{thm:p07-monotonicity-soundness}
 {proved=annex:p7.monotonicity.soundness}{Soundness}
Let \(S\) be a structure that satisfies every member of \(\Gamma\) under
every assignment.
Every assertion derivable from \(\Gamma\) is then true in \(S\) under
every assignment. Equivalently, derivability implies
\(\models_{\mathrm{mon}}\)-consequence.
\end{hbtheorem}

\begin{hbtheorem}{thm:p07-monotonicity-completeness}
 {proved=annex:p7.monotonicity.completeness}{Completeness with universal premises}
For every set of assertions \(\Gamma\),
\[
 \Gamma\models_{\mathrm{mon}}M\precsim_\tau N
 \quad\Longrightarrow\quad
 \Gamma\vdash_{\mathrm{mon}}M\precsim_\tau N .
\]
Together with soundness this gives equivalence. The models include
restricted function domains; completeness solely for full function
hierarchies is not asserted.
\end{hbtheorem}

\begin{hbtheorem}{thm:p07-monotonicity-lyndon}
 {proved=annex:p7.monotonicity.lyndon}
 {The polarity converse after coercion-aware elaboration}
Let \(\operatorname{tp}(M)\preceq\tau\) and \(x:\sigma_0\).
The following are equivalent:
\begin{enumerate}
\item Every free occurrence of \(x\) in \(L_\tau(M)\) is positive.
\item In every coherent structure and for every assignment \(\varphi\),
      the map
\[
 a\longmapsto
 \pi_{\operatorname{tp}(M),\tau}
      (\sem M_{\varphi_x^a}),\qquad a\in D_{\sigma_0},
\]
      is monotone.
\end{enumerate}
Replacing positive and monotone by negative and antitone gives the
corresponding equivalence. If \(x\) is absent from the long form, both
syntactic conditions hold and the displayed map is constant.
The criterion uses the specified comparison type and the receiving types
inside the elaboration. It is not a criterion on raw short normal syntax.
The model class includes restricted coherent structures.
\end{hbtheorem}

\begin{hbwarning}{def:p07-monotonicity-completion-target}
 {The scope of the adequacy and occurrence results}
The completeness and polarity-converse proofs use general coherent
structures, including syntactically constructed restricted function
domains. Neither result is a completeness claim solely about full
function hierarchies. The consequence relation has universally interpreted
premises throughout. The inference from raw unspecified polarity to
failure of semantic monotonicity is also unavailable, as
Theorem~\ref{thm:p07-monotonicity-coercion-polarity} demonstrates
\cite[Sections 8--10]{IcardMossTune2017Monotonicity}.
\end{hbwarning}

\section{Worked inference and failure cases}

\begin{hbexample}{def:p07-monotonicity-lexical}{Replacing a quantified term}
Interpret the predicate type \(p\) by subsets of a fixed universe, ordered
by inclusion, and truth values by \(0\leq1\).
Then
\[
 \mathsf{all}:p\xrightarrow-(p\xrightarrow+t),\qquad
 \mathsf{some}:p\xrightarrow+(p\xrightarrow+t).
\]
For example, \(\mathsf{all}(A)(B)=1\) precisely when \(A\subseteq B\).
If \(A\subseteq A'\), then every \(A'\) being a \(B\) entails every \(A\)
being a \(B\); if \(B\subseteq B'\), every \(A\) being a \(B\) entails
every \(A\) being a \(B'\). Nonempty intersection establishes the two
positive positions of \(\mathsf{some}\).

With \(\mathsf{poodle}\precsim_p\mathsf{dog}\), Anti at the first
argument of \(\mathsf{all}\), followed by Point, yields
\[
 \mathsf{all}(\mathsf{dog})(\mathsf{bark})
 \precsim_t
 \mathsf{all}(\mathsf{poodle})(\mathsf{bark}).
\]
Mono and Point give
\[
 \mathsf{some}(\mathsf{poodle})(\mathsf{bark})
 \precsim_t
 \mathsf{some}(\mathsf{dog})(\mathsf{bark}).
\]
The lexical inclusion is a premise, not a consequence of a word's spelling.
Likewise the set of plays lasting more than five hours is contained in the
set lasting more than two hours. Antitonicity of the subject position of
\(\mathsf{all}\) therefore turns \emph{Every play lasting more than two hours
is too long} into \emph{Every play lasting more than five hours is too long}.
The duration comparison supplies the set inclusion used by the rule.
\end{hbexample}

\begin{hbexample}{def:p07-monotonicity-neutral-failure}{An exact-count context}
Let \(\mathsf{two}(A)(B)\) assert that \(A\cap B\) has exactly two
elements. It is neither monotone nor antitone in either argument over all
sets. For a fixed set of barking dogs, enlarging a subject set from one to
two to three of them changes the truth values \(0,1,0\).
Consequently a \(\cdot\)-marked position licenses neither replacement
direction. In particular, two barking dogs that are not poodles make
\(\mathsf{two}(\mathsf{dog})(\mathsf{bark})\) true and
\(\mathsf{two}(\mathsf{poodle})(\mathsf{bark})\) false.
\end{hbexample}

\begin{hbexample}{def:p07-monotonicity-source-language-case}{The source's higher-order language example}
In addition to \(\mathsf{all}\), use
\(\mathsf{few}:p\xrightarrow{\cdot}(p\xrightarrow-t)\) and
\(\mathsf{deftly}:p\xrightarrow{\cdot}p\).
Suppose the following four assertions are supplied:
\[
 \begin{gathered}
 \lambda x.\mathsf{all}(x)(\mathsf{run})
   \precsim_{p\xrightarrow{\cdot}t}
 \lambda x.\mathsf{few}(x)(\mathsf{fly}),\\
 \mathsf{deftly}\precsim_{p\xrightarrow{\cdot}p}\lambda x.x,\qquad
 \mathsf{marsupial}\precsim_p\mathsf{mammal},\qquad
 \mathsf{soar}\precsim_p\mathsf{fly}.
 \end{gathered}
\]
Point and conversion give
\(\mathsf{deftly}(\mathsf{soar})\precsim_p\mathsf{soar}
 \precsim_p\mathsf{fly}\).
Anti for \(\mathsf{few}(\mathsf{marsupial})\) reverses this comparison.
The first abstraction on the first line is antitone in \(x\), so Anti
gives its mammal instance below its marsupial instance.
Point applied to the first premise then gives the corresponding
\(\mathsf{few}\) assertion. Transitivity combines these steps to yield
\[
 \mathsf{all}(\mathsf{mammal})(\mathsf{run})
 \precsim_t
 \mathsf{few}(\mathsf{marsupial})
                  (\mathsf{deftly}(\mathsf{soar})).
\]
The substantive bridge concerning running and flying is explicitly
assumed. The calculus does not establish it as a fact about animals.
\end{hbexample}

\begin{hbexample}{def:p07-monotonicity-source-arithmetic}{Instantiation in the source's arithmetic case}
Let \(2\) abbreviate \(+(1)(1)\). From \(0\precsim_r1\) and the
universally interpreted premise \(x\precsim_r +(x)(0)\), instantiate
\(x=1\). Mono for \(+(1)\) and Trans give \(1\precsim_r2\).
Mono in the first argument of subtraction gives
\(-(1)\precsim_{r\xrightarrow-r}-(2)\), hence Point gives
\(-(1)(1)\precsim_r-(2)(1)\).
Anti in the second argument and \(0\precsim_r1\) give
\(-(2)(1)\precsim_r-(2)(0)\). Thus
\(-(1)(1)\precsim_r-(2)(0)\), namely \(1-1\leq2-0\).

A separate analytical application uses the comparison theorem for
nonnegative improper integrals: on \(x\geq1\),
\(x\leq x^2\), so \(e^{-x^2}\leq e^{-x}\).
Together with nonnegativity and convergence of
\(\int_1^\infty e^{-x}\,dx=e^{-1}\), that theorem yields convergence of
\(\int_1^\infty e^{-x^2}\,dx\).
The calculus accounts for the variance steps; the integral comparison
theorem and its domain assumptions remain analytical premises.
\end{hbexample}

\begin{hbwarning}{def:p07-monotonicity-local-premises}{Why open premises are not local observations}
From a premise \(x\precsim_r0\) interpreted only at the assignment \(x=0\),
one cannot infer \(\lambda x.x\precsim_{r\xrightarrow{\cdot}r}
\lambda x.0\): evaluate both functions at \(1\).
Func is justified here because premises hold under every assignment.
The full-structure question and an alternative local-premise consequence
relation are different metatheoretic problems. This entry makes no
completeness claim for either alternative.
\end{hbwarning}

\begin{hbtheorem}{thm:p07-monotonicity-coercion-polarity}
 {proved=annex:p7.monotonicity.coercion-polarity}{A coercion can expose monotonicity}
Let \(b\) be a base type, \(x:b\) a variable, and put
\(\rho=b\xrightarrow+b\), \(\sigma=b\xrightarrow{\cdot}b\).
Let \(F:(\rho\xrightarrow+b)\xrightarrow+b\) be a constant.
The term
\[
 M=F(\lambda f:\sigma.\,f(x))
\]
is well typed, beta-eta normal, and its raw free occurrence of \(x\)
has polarity \(\cdot\). Nevertheless
\(a\mapsto\sem M_{\varphi_x^a}\) is monotone in every coherent
structure. The argument coercion restricts \(f\) to the monotone
function domain \(D_\rho\).
\end{hbtheorem}

\begin{hbmethod}{def:p07-monotonicity-lookup}{Using the entry}
First choose ordered semantic domains and type every constant.
Record lexical or mathematical inequalities as explicit premises.
Trace the polarity of the occurrence to be replaced, preserving argument
positions through application and abstraction. Use Mono in a positive
position and Anti in a negative one; a \(\cdot\) position needs further
information. Keep coercions and the comparison type visible whenever types
differ. Such coercions can expose additional signed behavior, as in
Theorem~\ref{thm:p07-monotonicity-coercion-polarity}; a raw unspecified
occurrence does not prove failure of semantic monotonicity.
Finally check whether a proposed premise is universal or only a
local observation, and state which nonlogical facts were assumed.
\end{hbmethod}
